\centerline{\bf \Huge Всеросс}

\begin{flushright}
        \scriptsize{Эпизод 7. Созданные человеком}
\end{flushright}


\problem{10.1}
Найдите количество корней уравнения

$$
|x|+|x+1|+\ldots+|x+2018|=x^2+2018 x-2019 .
$$


\problem{10.2}
Дан остроугольный треугольник $A B C$, в котором $A B<A C$. Пусть $M$ и $N$ - середины сторон $A B$ и $A C$ соответственно, а $D$ - основание высоты, проведённой из $A$. На отрезке $M N$ нашлась точка $K$ такая, что $B K=C K$. Луч $K D$ пересекает окружность $\Omega$, описанную около треугольника $A B C$, в точке $Q$. Докажите, что точки $C, N, K$ и $Q$ лежат на одной окружности.

\problem{10.3}
Дано натуральное число $k$. На клетчатой плоскости изначально отмечено $N$ клеток. Назовём крестом клетки $A$ множество всех клеток, находящихся в одной вертикали или горизонтали с $A$. Если в кресте неотмеченной клетки $A$ отмечено хотя бы $k$ других клеток, то клетку $A$ также можно отметить. Оказалось, что цепочкой таких действий можно отметить любую клетку плоскости. При каком наименьшем $N$ это могло случиться?

\problem{10.4}
Изначально на доске записано натуральное число. Затем каждую секунду к текущему числу прибавляют произведение всех его ненулевых цифр. Докажите, что найдётся натуральное $a$ такое, что прибавление числа $a$ случится бесконечное количество раз.